\section{Experimental Details and Additional Results}
\label{app:details}

\subsection{Source images and paired queries}
The retained GRScenes raster archive contains 80 paired views of 16 targets;
30 close-view pairs cover 12 of those targets. The 80 views include repeated observations of the same targets. The clean-control archive contains 15 pairs of nine
targets. Some early original-MDL captures contain material-resolution errors.
The later four-case material panel has repaired original captures and logged
provenance. We therefore distinguish archived raster readout experiments from
the clean-image window experiments used to probe material dependence. The archived readout analysis uses the saved raster inputs.

The two-view clean-image experiment uses the four repaired source pairs. It
forms four boards with distinct categories, two query targets per board and two
position permutations, yielding 16 paired queries per intervention. Source
IDs, image digests, boxes, and the exact composition recipe are recorded in
the manifest. The larger development composition uses 12 targets from the
archive and is recorded separately. All clean-image query and position variants reuse these four source pairs.

\subsection{Output interpretation}
In the text condition, the historical prompt requests \texttt{Point: x, y} and
\texttt{Answer: target category}. In the JSON condition it requests
\texttt{point\_xy} and \texttt{answer}. The target category is supplied in both
prompts, so answer-string agreement measures response adherence, not independent
object recognition. The new two-view task requests only a point. The action
variant names the target in an action description and tests action-framed object
localization, not fine-grained contact-point prediction.

For normalized frame size $K$, a generated point $(u,v)$ is mapped to image
pixels $(uW/K,vH/K)$ using the actual image dimensions. Pixel requests use the
point directly. Out-of-image coordinates and missing or unparsed points fail
the end-to-end metric. We preserve the raw response so that semantic errors,
format failures, and coordinate-scale mismatches can be examined separately.
The historical parser also handles four-number text lines by taking the box
center; this behavior is included in the parser tests and raw-output records.

\subsection{Archived readout experiment}
\label{app:archive}
Table~\ref{tab:archive} summarizes the common normalized-text readout. The
30-pair columns are exact-ID slices of the 80-pair inference runs, using the
same saved responses and current image digests. This avoids combining older
predictions with subsequently replaced images. The archive is an
appearance-stress input set for studying the readout interface; its render
issues are documented above.

\begin{table}[h]
\centering\small
\setlength{\tabcolsep}{3pt}
% Generated by scripts/build_assets.py from corrected results.json
\begin{tabular}{lrrrrrr}
\toprule
Model & \multicolumn{2}{c}{Close views (30)} & \multicolumn{2}{c}{Mixed views (80)} & $\Delta$ (pp) & Flips \\
\midrule
Gemma4-E4B & 90.0 & 96.7 & 70.0 & 65.0 & -5.0 & 12/80 \\
Qwen2.5-VL-7B & 10.0 & 13.3 & 3.8 & 5.0 & +1.2 & 5/80 \\
Qwen3-VL-2B & 86.7 & 86.7 & 53.8 & 50.0 & -3.8 & 13/80 \\
Qwen3-VL-8B & 93.3 & 96.7 & 86.2 & 85.0 & -1.2 & 7/80 \\
InternVL3-8B & 60.0 & 66.7 & 35.0 & 40.0 & +5.0 & 16/80 \\
Molmo-7B-D & 90.0 & 96.7 & 92.5 & 93.8 & +1.2 & 5/80 \\
\bottomrule
\end{tabular}

\caption{Archived raster readout: end-to-end point hits (\%), Original then
Converted within each pair of columns. $\Delta$ and flips refer to the 80-pair
pool. All queries remain in the denominator.}
\label{tab:archive}
\end{table}

\begin{table}[htbp]
\centering\small
\setlength{\tabcolsep}{3pt}
% Generated by scripts/build_assets.py from corrected results.json
\begin{tabular}{lllrrrr}
\toprule
Model & Frame & Format & $V_O$ & $V_C$ & $H_O$ (\%) & $H_C$ (\%) \\
\midrule
Gemma4-E4B & 0--1000 & text & 30/30 & 30/30 & 90.0 & 96.7 \\
Gemma4-E4B & 0--1000 & JSON & 30/30 & 30/30 & 90.0 & 86.7 \\
Gemma4-E4B & pixel & text & 30/30 & 30/30 & 36.7 & 36.7 \\
Gemma4-E4B & pixel & JSON & 30/30 & 30/30 & 0.0 & 0.0 \\
Qwen2.5-VL-7B & 0--1000 & text & 29/30 & 29/30 & 10.0 & 13.3 \\
Qwen2.5-VL-7B & 0--1000 & JSON & 9/30 & 12/30 & 13.3 & 6.7 \\
Qwen2.5-VL-7B & pixel & text & 26/30 & 25/30 & 80.0 & 66.7 \\
Qwen2.5-VL-7B & pixel & JSON & 18/30 & 22/30 & 60.0 & 73.3 \\
Qwen3-VL-2B & 0--1000 & text & 30/30 & 30/30 & 86.7 & 86.7 \\
Qwen3-VL-2B & 0--1000 & JSON & 30/30 & 30/30 & 90.0 & 93.3 \\
Qwen3-VL-2B & pixel & text & 30/30 & 30/30 & 60.0 & 60.0 \\
Qwen3-VL-2B & pixel & JSON & 30/30 & 30/30 & 0.0 & 0.0 \\
Qwen3-VL-8B & 0--1000 & text & 30/30 & 30/30 & 93.3 & 96.7 \\
Qwen3-VL-8B & 0--1000 & JSON & 30/30 & 30/30 & 93.3 & 96.7 \\
Qwen3-VL-8B & pixel & text & 30/30 & 30/30 & 0.0 & 3.3 \\
Qwen3-VL-8B & pixel & JSON & 30/30 & 30/30 & 0.0 & 0.0 \\
InternVL3-8B & 0--1000 & text & 25/30 & 26/30 & 60.0 & 66.7 \\
InternVL3-8B & 0--1000 & JSON & 11/30 & 7/30 & 36.7 & 20.0 \\
InternVL3-8B & pixel & text & 30/30 & 29/30 & 16.7 & 10.0 \\
InternVL3-8B & pixel & JSON & 11/30 & 16/30 & 10.0 & 20.0 \\
Molmo-7B-D & 0--1000 & text & 30/30 & 30/30 & 90.0 & 96.7 \\
Molmo-7B-D & 0--1000 & JSON & 30/30 & 30/30 & 13.3 & 3.3 \\
Molmo-7B-D & pixel & text & 30/30 & 30/30 & 0.0 & 0.0 \\
Molmo-7B-D & pixel & JSON & 30/30 & 30/30 & 0.0 & 0.0 \\
\bottomrule
\end{tabular}

\caption{Full readout matrix. $V$ is the number of valid numeric points over all
30 queries; $H$ is end-to-end point-hit rate. A valid numeric point can still be
out of frame or outside the target. The two-line normalized-text rows are
exact-ID subsets of current 80-pair runs; other rows use their recorded 30-pair
inference runs.}
\label{tab:protocol-full}
\end{table}

\subsection{Window interventions and uncertainty}
The target-window crossover uses the floor/ceiling integer bounds of the
projected target box. Pixels inside that rectangle are taken from one material
condition and pixels outside from the other. The rectangle includes background
and is not an instance segmentation mask. The mechanism experiment identifies
image-region dependence; it does not individually isolate shader parameters
such as roughness or opacity. The clean source set contains bottle, clock,
cup, and a bag annotated as backpack. The annotation names are retained in the
queries. 

We report fixed-denominator paired differences. Target-cluster bootstrap
resamples all views of an object together; the two-view analysis groups the
query and position permutations by board. Boards share source targets, so this
bootstrap describes board variability rather than a new-object population.
Repeated trials are never counted as additional independent targets. Medians
of even-sized samples average the two middle observations. Pixel displacements
are computed only for pairs with two numeric points, with that count retained
separately from the success-rate denominator.

\subsection{Model environments and decoding}
Historical Qwen/InternVL runs use Transformers 5.2 and PyTorch 2.10. The
Gemma4 four-bit checkpoint uses its Unsloth-compatible environment. Molmo uses
a dedicated Transformers 4.46.3/PyTorch 2.5.1 environment because its remote
generation code expects a legacy cache. InternVL uses the remote chat model
with a single 448-pixel square image tile. These are model-plus-preprocessor
configurations; their differing image pipelines are recorded in the artifact.
The primary runs use greedy decoding. The historical grid uses a 64-token
output budget; the new two-view probes use 128 tokens. Checkpoint revisions,
runtime patches, prompts and settings are retained with each run.

\subsection{Stochastic spatial readout}
The recorded 30-pair raster inputs are evaluated at temperature 0.7 with three
seeds per model. For each image, seed distance is the mean of the three
pairwise pixel distances. A query has variable hit status if its three
end-to-end outcomes are not identical. Missing points are failures in this
last metric. All 60 queries in each of the three model groups have three
numeric points. These runs measure decoding variability on the archived
inputs, separately from the clean-image material-region experiment.

\begin{table}[htbp]
\centering\small
% Generated by scripts/build_assets.py from corrected results.json
\begin{tabular}{lrrr}
\toprule
Model & Complete triples & Seed distance (px) & Hit variability (\%) \\
\midrule
Qwen3-VL-8B & 60/60 & 7.66 & 0.0 \\
Qwen3-VL-2B & 60/60 & 24.34 & 16.7 \\
Gemma4-E4B & 60/60 & 27.07 & 10.0 \\
\bottomrule
\end{tabular}

\caption{Stochastic readout on identical images, recomputed from the saved
three-seed predictions. Complete triples count queries with three numeric points.}
\label{tab:sampling}
\end{table}

\subsection{Navigation pairing and recorded examples}
The retained KuJiaLe package contains scenes 0031, 0036 and 0066, each with
33 episodes. The analysis retains the full set of 99 paired IDs. Each episode has one metric rollout per material condition, so the observed
turnover includes any runtime variability; there is no same-condition repeat
estimate to subtract. GRScenes grounding and KuJiaLe navigation use different
scenes and models and provide complementary levels of observation, without
establishing a causal chain between a grounding error and a navigation outcome.

Full navigation metric runs are video-free. Selected visual examples come from
separate recorded reruns with corresponding episode IDs. Their media illustrate
trajectory behavior; their scores do not replace the full-run results.

\subsection{Assets and release}
The experimental sources include third-party scenes and model checkpoints.
The local submission-preparation package contains the manuscript, analysis
tables, and reproduction instructions. Scene redistribution follows the
respective source terms. Private machine paths and author-identifying metadata
are excluded from the anonymous PDF and its copy-ready submission text.
